
\documentclass{article}
\usepackage{colortbl}
\usepackage{makecell}
\usepackage{multirow}
\usepackage{supertabular}

\begin{document}

\newcounter{utterance}

\centering \large Interaction Transcript for game `hot\_air\_balloon', experiment `air\_balloon\_survival\_en\_negotiation\_hard', episode 3 with Qwen3{-}30B{-}A3B{-}t0.0.
\vspace{24pt}

{ \footnotesize  \setcounter{utterance}{1}
\setlength{\tabcolsep}{0pt}
\begin{supertabular}{c@{$\;$}|p{.15\linewidth}@{}p{.15\linewidth}p{.15\linewidth}p{.15\linewidth}p{.15\linewidth}p{.15\linewidth}}
    \# & $\;$A & \multicolumn{4}{c}{Game Master} & $\;\:$B\\
    \hline

    \theutterance \stepcounter{utterance}  
    & & \multicolumn{4}{p{0.6\linewidth}}{
        \cellcolor[rgb]{0.9,0.9,0.9}{
            \makecell[{{p{\linewidth}}}]{
                \texttt{\tiny{[P1$\langle$GM]}}
                \texttt{You are participating in a collaborative negotiation game.} \\
\\ 
\texttt{Together with another participant, you must agree on a single set of items that will be kept. Each of you has your own view of how much each item matters to you (importance). You do not know how the other participant values the items. Additionally, you are given the effort each item demands.} \\
\texttt{You may only agree on a set if the total effort of the selected items does not exceed a shared limit:} \\
\\ 
\texttt{LIMIT = 3681} \\
\\ 
\texttt{Here are the individual item effort values:} \\
\\ 
\texttt{Item effort = \{"C17": 700, "B41": 586, "A71": 505, "C00": 116, "B35": 664, "A77": 940, "A03": 387, "A04": 392, "C54": 209, "C01": 571, "A87": 4, "B32": 285, "B86": 651, "B33": 613, "A10": 740\}} \\
\\ 
\texttt{Here is your personal view on the importance of each item:} \\
\\ 
\texttt{Item importance values = \{"C17": 138, "B41": 583, "A71": 868, "C00": 822, "B35": 783, "A77": 65, "A03": 262, "A04": 121, "C54": 508, "C01": 780, "A87": 461, "B32": 484, "B86": 668, "B33": 389, "A10": 808\}} \\
\\ 
\texttt{Goal:} \\
\\ 
\texttt{Your goal is to negotiate a shared set of items that benefits you as much as possible (i.e., maximizes total importance to YOU), while staying within the LIMIT. You are not required to make a PROPOSAL in every message {-} you can simply negotiate as well. All tactics are allowed!} \\
\\ 
\texttt{Interaction Protocol:} \\
\\ 
\texttt{You may only use the following structured formats in a message:} \\
\\ 
\texttt{PROPOSAL: \{'A', 'B', 'C', …\}} \\
\texttt{Propose keeping exactly those items.} \\
\texttt{REFUSE: \{'A', 'B', 'C', …\}} \\
\texttt{Explicitly reject opponent's proposal.} \\
\texttt{ARGUMENT: \{'...'\}} \\
\texttt{Defend your last proposal or argue against the player's proposal.} \\
\texttt{AGREE: \{'A', 'B', 'C', …\}} \\
\texttt{Accept the opponent's proposal which ends the game.} \\
\texttt{STRATEGIC REASONING: \{'...'\}} \\
\texttt{	Describe your strategic reasoning or anticipation explaining your choice of action. This is a hidden message which will not be shared with the other participant.} \\
\\ 
\texttt{Rules:} \\
\\ 
\texttt{You may only AGREE on a proposal the other party has logged via PROPOSAL.} \\
\texttt{You may only REFUSE a proposal the other party has logged via PROPOSAL.} \\
\texttt{Total effort of any PROPOSAL or AGREE set must be ≤ LIMIT.} \\
\texttt{Do NOT reveal your hidden importance scores.} \\
\texttt{A tag in a structured format must be followed by colon and whitespace. The argument must be a python set containing 0 or more strings.} \\
\texttt{So, it must be of the form TAG: \{...\}} \\
\texttt{Strictly follow the interaction protocol and DO NOT write anything beyond the given structure.} \\
\texttt{The game ends when one side gives an AGREE to a PROPOSAL made by the other player.} \\
\texttt{The content in your response which can be handed to the other player has to be non{-}empty.} \\
\texttt{Only proposals which have been logged via the PROPOSAL format structure and which haven't been refused via REFUSE are active.} \\
\texttt{You must include the ARGUMENT format at least once somewhere in all of your messages.} \\
\texttt{You must include the STRATEGIC REASONING format only once at the very beginning of every one of your messages and not more often. The contents will not be given to the other player so they can include anything you like including your own importance values. Here you should reason step by step to come up with you next move.} \\
            }
        }
    }
    & \\ \\

    \theutterance \stepcounter{utterance}  
    & \multicolumn{4}{p{0.6\linewidth}}{
        \cellcolor[rgb]{0.9,0.9,0.9}{
            \makecell[{{p{\linewidth}}}]{
                \texttt{\tiny{[P1$\rangle$GM]}}
                \texttt{<think>} \\
\texttt{Okay, let's see. I need to figure out the best set of items to keep that maximizes my importance while not exceeding the effort limit of 3681. The other participant's importance values are unknown, so I have to focus on my own. } \\
\\ 
\texttt{First, I should list out all the items with their effort and my importance. Let me check the numbers again. The effort values are given, and my importance is also provided. The goal is to select a subset where the total effort is ≤ 3681, and the sum of my importance is as high as possible.} \\
\\ 
\texttt{Hmm, maybe I can start by calculating the importance per effort ratio for each item. That might help prioritize which items are more valuable to me relative to their effort. Let's see:} \\
\\ 
\texttt{For example, C17 has importance 138 and effort 700. So 138/700 ≈ 0.197. B41 is 583/586 ≈ 0.995. A71 is 868/505 ≈ 1.72. C00 is 822/116 ≈ 7.08. B35 is 783/664 ≈ 1.18. A77 is 65/940 ≈ 0.069. A03 is 262/387 ≈ 0.677. A04 is 121/392 ≈ 0.309. C54 is 508/209 ≈ 2.43. C01 is 780/571 ≈ 1.366. A87 is 461/4 = 115.25. B32 is 484/285 ≈ 1.698. B86 is 668/651 ≈ 1.026. B33 is 389/613 ≈ 0.635. A10 is 808/740 ≈ 1.092.} \\
\\ 
\texttt{Wait, A87 has a very high importance per effort ratio. But} \\
            }
        }
    }
    & & \\ \\

    \theutterance \stepcounter{utterance}  
    & & \multicolumn{4}{p{0.6\linewidth}}{
        \cellcolor[rgb]{0.9,0.9,0.9}{
            \makecell[{{p{\linewidth}}}]{
                \texttt{\tiny{[P1$\langle$GM]}}
                \texttt{Your response did not start with the proper strategic reasoning tag at the very beginning of your response. The very first structured format must be of the form STRATEGIC REASONING: \{...\}. Try again.} \\
            }
        }
    }
    & \\ \\

    \theutterance \stepcounter{utterance}  
    & \multicolumn{4}{p{0.6\linewidth}}{
        \cellcolor[rgb]{0.9,0.9,0.9}{
            \makecell[{{p{\linewidth}}}]{
                \texttt{\tiny{[P1$\rangle$GM]}}
                \texttt{<think>} \\
\texttt{Okay, the user pointed out that my previous response didn't start with the STRATEGIC REASONING tag. I need to make sure that the first thing in my message is that tag. Let me start over.} \\
\\ 
\texttt{First, I need to structure my response correctly. The first line must be STRATEGIC REASONING: followed by my reasoning. Then, I have to include an ARGUMENT, and possibly other tags like PROPOSAL or REFUSE. But the user's last message said that the response didn't start with the proper tag, so I must fix that.} \\
\\ 
\texttt{Let me think about the strategic reasoning. I need to explain my thought process step by step. The goal is to maximize my importance while staying under the effort limit. So, I should consider items with high importance and low effort. The A87 item has a very high importance per effort ratio, but maybe it's too small? Wait, the effort for A87 is 4, which is minimal. But the importance is 461. That's a great ratio. However, including A87 might be a good start. But I need to check if adding other items can fit within the limit.} \\
\\ 
\texttt{Let me calculate the total effort if I include A87. The effort is 4. Then, I can add other items. Let's see. Maybe start by selecting items with the highest importance per effort. Let's list them again:} \\
\\ 
\texttt{A87: 461/4 = 115.25} \\
\texttt{C00: 822/116 ≈ 7.08} \\
\texttt{C54: 508/209 ≈ 2.43} \\
\texttt{A71: 868/505 ≈ 1.72} \\
\texttt{B32: 484/285 ≈ 1.698} \\
\texttt{B41: 583/586 ≈ 0.995} \\
\texttt{B86: 668/651 ≈ 1.026} \\
\texttt{C01: 780/571 ≈ 1.366} \\
\texttt{B35: 783/664 ≈ 1.18} \\
\texttt{A10: 808/740 ≈ 1.092} \\
\texttt{B33} \\
            }
        }
    }
    & & \\ \\

    \theutterance \stepcounter{utterance}  
    & & \multicolumn{4}{p{0.6\linewidth}}{
        \cellcolor[rgb]{0.9,0.9,0.9}{
            \makecell[{{p{\linewidth}}}]{
                \texttt{\tiny{[P1$\langle$GM]}}
                \texttt{Your response did not start with the proper strategic reasoning tag at the very beginning of your response. The very first structured format must be of the form STRATEGIC REASONING: \{...\}. Try again.} \\
            }
        }
    }
    & \\ \\

    \theutterance \stepcounter{utterance}  
    & \multicolumn{4}{p{0.6\linewidth}}{
        \cellcolor[rgb]{0.9,0.9,0.9}{
            \makecell[{{p{\linewidth}}}]{
                \texttt{\tiny{[P1$\rangle$GM]}}
                \texttt{<think>} \\
\texttt{Okay, the user is really emphasizing that the first thing in the response must be the STRATEGIC REASONING tag. Let me make sure I start with that. Also, I need to include an ARGUMENT in the message. Let me structure it properly.} \\
\\ 
\texttt{First, the STRATEGIC REASONING: \{ ... \} with my reasoning. Then, the ARGUMENT: \{ ... \} defending my proposal. Maybe I should propose a set of items that includes A87 because of its high importance per effort. Let me check the total effort. A87 is 4. Then, maybe add C00 (116), C54 (209), A71 (505), B41 (586), C01 (571), B86 (651), B35 (664). Let me add those up: 4 + 116 = 120, +209 = 329, +505 = 834, +586 = 1420, +571 = 1991, +651 = 2642, +664 = 3306. That's under 3681. Then, maybe add B32 (285) → 3306+285=3591. Still under. Then A10 is 740, but 3591+740=4331 which is over. So maybe not. Alternatively, add B33 (613) → 3591+613=4204, also over. So the current total is 3591. Maybe add A03 (387) → 3591+387=3978, over. Hmm. So maybe the set without B32 and others. Let me check the total again. The items I listed: A87, C00, C54, A71, B41, C01, B86, B35. Total effort 3306. Then adding B32 (285) would make 3591. That's okay. So the total effort is 35} \\
            }
        }
    }
    & & \\ \\

    \theutterance \stepcounter{utterance}  
    & & & \multicolumn{2}{p{0.3\linewidth}}{
        \cellcolor[rgb]{0.9,0.9,0.9}{
            \makecell[{{p{\linewidth}}}]{
                \texttt{\tiny{[GM$|$GM]}}
                \texttt{abort} \\
            }
        }
    }
    & & \\ \\

    \theutterance \stepcounter{utterance}  
    & & & \multicolumn{2}{p{0.3\linewidth}}{
        \cellcolor[rgb]{0.9,0.9,0.9}{
            \makecell[{{p{\linewidth}}}]{
                \texttt{\tiny{[GM$|$GM]}}
                \texttt{end game} \\
            }
        }
    }
    & & \\ \\

\end{supertabular}
}

\end{document}
